% Introduction. Sources: results/probes/binary_*_test.md,
% results/chartqa_type_probes.md, results/labels/*.error_stats.json,
% results/synth_fabrication_pilots.md, results/pilot3_probes.md, results/e4.md,
% results/e5.md.
\section{Introduction}
\label{sec:intro}

Vision-language models (VLMs) are routinely asked to read charts, and a wrong
answer can be wrong in two different ways. The model may have \emph{misread}
the figure: taken the adjacent bar, rounded to a gridline, or misread a printed
value. The answer was in the image, and a closer look might recover it. Or it
may have \emph{fabricated}: reported a value the figure does not contain.
Looking again cannot help, and the right response is to withhold the answer.
Reliability methods usually collapse both cases into one ``do not trust''
score. We ask whether the two failure types are visible in the model's internal
state \emph{before} it answers, and whether knowing the type helps choose a
repair.

\paragraph{Problem.} Let $x=(I,q)$ be a chart image and a question, and $\hat y$
the answer a frozen VLM generates greedily. Let $h_\ell(x)\in\mathbb{R}^d$ be the
hidden state after decoder layer $\ell$ at the last prompt token, computed in the
same forward pass that produces $\hat y$. Each wrong answer has a failure type
$t$, read off $\hat y$ and the figure. We test (i) whether linear probes on
$h_\ell$ predict $\Pr(\hat y \text{ wrong})$ and $t$ better than surface features
of the input, (ii) whether a probe trained for one type also responds to the
other (cross-transfer), (iii) whether the best repair depends on $t$ (E4), and
(iv) whether a controller that uses the type beats plain selective prediction
(E5). The \textbf{Type Separability Hypothesis} states that the states before a
misread and before a fabrication are linearly distinguishable, the first being
a precision-limited perceptual failure and the second an answer ungrounded in
the image. A collapse of the two classes is a reportable outcome, not a failure
of the design.

\paragraph{Contributions.} (1) Full-depth activation caches of
\qwen and \llava on ChartQA and on three synthetic pilots, taken from the pass
that produced each answer. (2) Five-label annotations of all 1,495 ChartQA
errors of both models, with a two-person agreement check. (3) Probes that
predict wrong answers (AUROC 0.877 and 0.798) and separate misreads from
arithmetic errors on real charts (0.897 and 0.961), each against surface
baselines, residualisation, five data splits and single-question-type
controls. (4) A synthetic arm that induces fabrication: neither model refuses
questions about a missing bar, yet both encode that it is missing (AUROC 0.96 to
1.00), and the same signal fires on years outside a chart's range. (5) A
pre-registered repair experiment in which cropping rescues misreads more than
fabrications in both models, and a controller evaluated against selective
prediction.
